Let the outer side be $1$ and diagonal be $\phi$. Similar triangles in the pentagram give $\phi=1+\dfrac{1}{\phi}$, so $\phi=\dfrac{1+\sqrt5}{2}$. The central side is the diagonal with two end pieces of length $\dfrac{1}{\phi}$ removed, hence $\phi-\dfrac{2}{\phi}=2-\phi=\dfrac{1}{\phi^2}$. The two pentagons are similar, so their area ratio is $\dfrac{1}{\phi^4}=\dfrac{2}{7+3\sqrt5}$. Thus $(a,b,c,d)=(2,7,3,5)$ and the sum is $\boxed{17}$.
\begin{asy}
texpreamble("\usepackage{amsmath,amssymb}");
size(220);
defaultpen(fontsize(9pt));
pair[] P;for(int i=0;i<5;++i)P.push(dir(90+72*i));for(int i=0;i<5;++i){draw(P[i]--P[(i+1)%5]);draw(P[i]--P[(i+2)%5]);label("$"+string(i+1)+"$",P[i],P[i]);}
\end{asy}
